\subsection{Controller Workflow}
\label{sec:workflow}

Algorithm~\ref{alg:overview} summarizes how the controller in Figure~\ref{fig:controller-overview} coordinates admission and pacing to produce draft images \(\mathit{draft}_1,\ldots,\mathit{draft}_N\) from shutter inputs for captures \(1,\ldots,N\).
Let \(\mathcal K=(k_1,\ldots,k_m)\) denote the configured draft sequence for consecutive captures.
Each stage is represented by a workload key \(k\) that identifies both the stage and its configuration.
The final key \(k_m\) represents mandatory encoding and saving.
To support admission and pacing decisions, the algorithm initializes an empty duration history \(\mathcal D\) to store measured processing times from completed draft sequences.
For pacing, it also uses \(\mathit{deadline}_{\mathrm{queue}}\), initially set to zero, to record the deadline of the most recently enqueued capture in the draft queue (line~\ref{alg:ov:configure}).

\begin{algorithm}[t]
\caption{Capture Workflow with Admission and Pacing}
\label{alg:overview}
\footnotesize
\renewcommand{\algorithmicrequire}{\textbf{Input:}}
\renewcommand{\algorithmicensure}{\textbf{Output:}}
\algdef{SE}[DRAFTASYNC]{AsyncDraft}{EndAsyncDraft}
  {\textbf{Asynchronously do:}}
  {\algorithmicend}
\algtext*{EndIf}\algtext*{EndFor}\algtext*{EndAsyncDraft}
\begin{algorithmic}[1]
\Require Shutter inputs for captures $1,\ldots,N$
\Ensure Draft images $\mathit{draft}_1,\ldots,\mathit{draft}_N$
\State $\mathcal K_i\gets(k_1,\ldots,k_m)\;(\forall i),\quad\mathcal D\gets\emptyset,\quad\mathit{deadline}_{\mathrm{queue}}\gets0$ \label{alg:ov:configure}\label{alg:ov:history}\label{alg:ov:queue-init}
\State
\For{$i=1$ \textbf{to} $N$}
  \State $\mathit{deadline}_i\gets\textsf{Now}()+\eta_{\mathrm{timeout}}$ \Comment{$\eta_{\mathrm{timeout}}=$ fixed timeout duration} \label{alg:ov:timeout}
  \State $T_{\mathrm{queue}}\gets\mathit{deadline}_{\mathrm{queue}}-\textsf{Now}()$ \label{alg:ov:queue-budget}
  \State $\mathit{delay}_i\gets\Call{DecidePacingDelay}{\mathcal D,T_{\mathrm{queue}}}$ \label{alg:ov:pacing}
  \State
  \State $\mathit{deadline}_{\mathrm{queue}}\gets\mathit{deadline}_i$ \label{alg:ov:queue-update}
  \AsyncDraft \label{alg:ov:async}
    \For{each optional stage $k_j\in\mathcal K_i$}
      \State $\mathcal K_{i,j}\gets(k_j,\ldots,k_m)$ \label{alg:ov:remaining}
      \State $T_{i,j}\gets\mathit{deadline}_i-\textsf{Now}()$
      \If{$\Call{DecideAdmission}{\mathcal D,\mathcal K_{i,j},T_{i,j}}=\textsc{admit}$} \label{alg:ov:admission}
        \State $\textsf{Execute}(k_j)$
      \EndIf
    \EndFor
    \State $(\mathit{draft}_i,p_i,C_i)\gets\textsf{EncodeAndSave}(i)$ \label{alg:ov:publish}
    \State $\mathcal D\gets\Call{Update}{\mathcal D\cup\{(\mathcal K_i,p_i,C_i)\}}$ \label{alg:ov:record}\label{alg:ov:update}
  \EndAsyncDraft
  \State $\textsf{Wait}(\mathit{delay}_i)$ \label{alg:ov:wait}
\EndFor
\State \Return $(\mathit{draft}_1,\ldots,\mathit{draft}_N)$
\end{algorithmic}
\end{algorithm}


For each accepted capture \(i\), the algorithm sets \(\mathit{deadline}_i\) using the fixed timeout duration \(\eta_{\mathrm{timeout}}\) (line~\ref{alg:ov:timeout}).
When \texttt{captureAvailable} is received from the HAL, the controller first computes \(T_{\mathrm{queue}}\), the time remaining before \(\mathit{deadline}_{\mathrm{queue}}\) (line~\ref{alg:ov:queue-budget}).
It then uses \textsc{DecidePacingDelay} to compute \(\mathit{delay}_i\) by comparing draft workload estimated from the duration history \(\mathcal D\) with \(T_{\mathrm{queue}}\) (line~\ref{alg:ov:pacing}).
In the original workflow, the signal is delivered immediately (\(\mathit{delay}_i=0\)), whereas the controller schedules its delivery using the computed delay, determining when capture \(i+1\) can be accepted (line~\ref{alg:ov:wait}).

When capture \(i\)'s draft sequence enters the draft queue, the algorithm sets \(\mathit{deadline}_{\mathrm{queue}}\) to \(\mathit{deadline}_i\), updating the deadline used in subsequent pacing decisions (line~\ref{alg:ov:queue-update}).
Capture \(i\)'s draft processing runs asynchronously on a single draft worker (line~\ref{alg:ov:async}), which repeatedly executes queued sequences in capture order until the draft queue drains.
This processing proceeds independently of the callback delivery delay and can overlap acceptance and frame collection for subsequent captures.
The controller uses \textsc{DecideAdmission} to decide whether optional stages execute at runtime.
The existing static thermal guard excludes optional stages during draft sequence configuration, before processing begins (Section~\ref{sec:guard-limit}), whereas the controller retains these stages in \(\mathcal K\) for admission.
Before each optional stage \(k_j\), the controller obtains the remaining sequence \(\mathcal K_{i,j}=(k_j,\ldots,k_m)\) (line~\ref{alg:ov:remaining}) and the capture's remaining timeout budget \(T_{i,j}\) (line~\ref{alg:ov:remaining-budget}).
It then uses the duration history \(\mathcal D\) to compute an empirical upper estimate of the remaining processing time, including mandatory encoding and saving, and compares this estimate with \(T_{i,j}\) to decide whether \(k_j\) executes (line~\ref{alg:ov:admission}).

To complete the draft sequence, the draft worker performs mandatory encoding and saving through \textsf{EncodeAndSave}, which returns \(\mathit{draft}_i\) and the collected duration measurements (line~\ref{alg:ov:publish}).
These measurements include the stage durations \(p_i\), with \(p_i(k)\) denoting the processing duration of each executed stage \(k\), and the draft sequence duration \(C_i\).
The controller records these measurements with \(\mathcal K\) in \(\mathcal D\) and updates the shared duration predictor and the admission and pacing states through \textsc{Update} for subsequent decisions (line~\ref{alg:ov:update}).

